\uuid{BU6H}
\exo7id{1576}
\titre{exo7 1576}
\auteur{legall}
\organisation{exo7}
\datecreate{1998-09-01}
\isIndication{false}
\isCorrection{false}
\chapitre{Réduction d'endomorphisme, polynôme annulateur}
\sousChapitre{Polynôme annulateur}
\cpge{mpsi-mp2i.matrices-applications-lineaires.a, pcsi.matrices-applications-lineaires.a, ptsi.matrices-applications-lineaires.a, mpsi-mp2i.changements-bases.a, pcsi.matrices-applications-lineaires.c, ptsi.matrices-applications-lineaires.c, mp-mpi.reduction-endomorphismes.b, pc.reduction.a, psi.reduction-endomorphismes.a, pt.reduction-endomorphismes.a, mp-mpi.reduction-endomorphismes.d, pc.reduction.c, psi.reduction-endomorphismes.c, pt.reduction-endomorphismes.c, mp-mpi.reduction-endomorphismes.g, pc.complements-espaces-vectoriels.d, psi.complements-espaces-vectoriels.d, mp-mpi.reduction-endomorphismes.i, pc.reduction.d, psi.reduction-endomorphismes.d, bcpst2.valeurs-propres-vecteurs-propres.a, bcpst2.valeurs-propres-vecteurs-propres.b}
\module{Algèbre}
\niveau{L2}
\difficulte{}

\contenu{
\texte{
Soit $  E  $ un $  \mathbb{K} $-espace vectoriel de dimension finie
$  n  $ et $  f
\in \mathcal{L} (E)  $ tel que $  \hbox{rg}(f-id)=1  .$ On note $  H = \hbox{Ker}(f-id)  .$
}
\begin{enumerate}
    \item \question{Soit $  \{ e_1, \cdots , e_{n-1}\}   $ une base de $  H   $ et $  e_n \notin H  .$ Montrer que $  \{ e_1, \ldots , e_{n}\}   $ est une base de $  E   $ et donner l'allure de la matrice de $  f   $ dans cette base.}
    \item \question{Montrer que le polyn\^ome $  (X-1)(X-\hbox{det}(f))  $ annule $  f  .$ Donner une condition n\'ec\'essaire et
suffisante pour que $  f  $ soit diagonalisable.}
\end{enumerate}
}
